% Introduction (Task 46, Scientific Reports): v3 Introduction and Related Work condensed (plan 42, section 1.3); no
% subheadings. Every result number is a macro.
\section*{Introduction}
\label{sec:intro}

Machine-learning network intrusion detection systems (NIDS) classify network flows from statistical features, and
transformer models for tabular data have been applied to this task \cite{gorishniy2021,manocchio2024,borisov2024}.
Their attention weights are an attractive explanation because they come with every prediction at no extra cost, and
tabular transformers have used attention to describe feature importance \cite{arik2021,gorishniy2021,somepalli2021}.
Whether attention is a faithful explanation, one that reflects what drives the prediction, has been argued in both
directions for text models \cite{jain2019,wiegreffe2019,serrano2019,jacovi2020}, and methods that aggregate attention
across layers or combine it with gradients were proposed because raw attention weights mix information across layers
\cite{abnar2020,chefer2021,ferrando2022}.
For flow-feature models the question has received less attention.
Evaluations of explanation methods in security have proposed criteria for this setting \cite{warnecke2020}, shown
that SHAP explanations of NIDS depend on the replacement data \cite{tritscher2023}, and compared white-box methods on
intrusion datasets \cite{arreche2025}.
Explanations have also been built into ensemble intrusion detectors, where SHAP-based explanations describe the
detector or select its features \cite{ahmed2024haenid,ahmed2024shap,ahmed2025consensus,ahmed2025hcivad}.

Two further questions arise once an attacker is present.
First, an attacker who perturbs a flow to evade the classifier may also change the explanation.
Explanations can be changed by small input perturbations without changing the prediction
\cite{ghorbani2019,dombrowski2019}, attention can be made deceptive through training \cite{pruthi2020}, and
explanations can be manipulated through the model \cite{heo2019,slack2020}; these attacks and their countermeasures
have been surveyed \cite{baniecki2024}.
How attention behaves under attack therefore matters for any analyst who reads it.
Second, if attention and attribution methods agree on clean traffic but disagree on adversarial traffic, their
consistency could serve as an attack signal.
Explanation-based detection of adversarial inputs has been proposed with SHAP signatures \cite{fidel2020}, attacks that
mislead a classifier and its interpretation together have been built \cite{zhang2020}, and for intrusion detection the
drift of SHAP explanations under gradient attacks has been used as an auxiliary detection signal \cite{maseno2026}.
Detectors of adversarial inputs have repeatedly been bypassed by attacks designed for them, so such a detector has to
be evaluated against attackers that model it \cite{carlini2017,tramer2020}, and the strength of that attacker decides
what the evaluation can show.

Both questions only make sense for attacks whose flows respect the protocol.
Feature-space attacks on NIDS can produce flows that no real traffic could generate, which motivates problem-space
constraints, validity checks and realistic threat models \cite{pierazzi2020,alhussien2024,apruzzese2022}; constrained
attacks on NIDS include saliency-map attacks under domain constraints \cite{sheatsley2022} and black-box botnet evasion
\cite{debicha2023}, and learning-based security evaluations have known pitfalls \cite{arp2022}.
We therefore restricted every attacker to a hand-built model of flow-feature constraints and checked every adversarial
row with a separate validity checker.

We asked three research questions.
Is attention a faithful explanation of a flow classifier, on clean flows and under protocol-constrained attacks (RQ1)?
Do attention-attribution consistency scores improve a confidence-and-distance attack detector, first against an
attacker that ignores them and then against attackers that model them, from differentiable stand-ins to the
detector's exact score, and is any gain specific to attention (RQ2)?
How much do the answers depend on the dataset and its data quality (RQ3)?
We answered them with a feature-tokenizer transformer (FT-Transformer) trained with \nSeeds{} seeds on two corrected
intrusion datasets, LYCOS-IDS2017 and LycoS-Unicas-IDS2018, using the \nCommonClasses{} attack classes common to both
(Fig.~\mainFigPipeline).
The attack detector was evaluated against three levels of attacker: one that ignores the consistency scores, one that
models them through differentiable stand-ins, and one that optimises the exact detector score with surrogate gradients
only for the rank-based scores, compared under an equal-effort protocol with outcome rules recorded before the runs.
\nControlsAll{} control detectors replace attention by other inputs to separate its role.

Our central finding concerns the second question.
Against attackers that optimised the exact score of the detector they faced, with surrogate gradients only for the
\nRankScores{} rank-based consistency scores, adding the consistency scores raised the share of successful and
undetected attack rows, by \axHeadVoneSeventeenAbs{} on 2017 and by \axHeadVoneEighteenAbs{} on 2018, with
\axHeadOppSeedsPhrase{} of the opposite sign.
An attacker with differentiable stand-ins for the same scores had shown a reduction of \axStandinVoneSeventeenAbs{}
on 2017, so that result did not hold against the stronger attacker.
Giving the attacker of the base detector the same iterate selection as the stronger attacker changed the outcome of
\sxOutcomeChangedPhrase, so the reversal does not come from the selection rule.
For RQ1, attention changed little under constrained attacks but was less faithful than integrated gradients and
occlusion.

The study makes three contributions.
First, it compares three attention variants with four attribution methods on two intrusion datasets, by agreement and
by deletion and insertion faithfulness on clean data and on protocol-constrained adversarial rows, and shows that the
loss of faithfulness after an evasion follows the prediction flip rather than the explainer.
Second, it evaluates consistency monitoring against three levels of attacker under an equal-effort protocol with outcome
rules recorded before the runs, and gives a concrete case in which an adaptive attack with differentiable stand-ins
overstated what a detector achieves.
Third, it separates the role of attention with control detectors, \nControlsAttacked{} of them attacked by adversaries
that modelled their own scores; the gain against the ignoring attacker was largely obtained without attention, and no
control was clearly better than the base detector on both datasets under its own adaptive attacker.
The constraint model with its validity checker and the near-duplicate analysis support these results.
We do not propose a new explanation method or a new detector, and we do not claim that attention is unfaithful in
general; the results refer to one model family and two datasets.
